\documentclass[aspectratio=169]{beamer}
\usetheme{Madrid}
\usecolortheme{seahorse}

% Packages for symbols and tables
\usepackage{pifont}
\newcommand{\done}{\textcolor{green}{\ding{51}}} % ✓
\newcommand{\inprogress}{\textcolor{orange}{\ding{217}}} % →
\newcommand{\pending}{\textcolor{red}{\ldots}} % …

% Custom footer: PBADT text left + date center + page numbers right
\setbeamertemplate{footline}{%
  \leavevmode%
  \hbox{%
    \begin{beamercolorbox}[wd=.33\paperwidth,ht=2.5ex,dp=1ex,left]{author in head/foot}%
      \hspace*{2mm}\scriptsize PBADT504 BIG DATA ANALYTICS
    \end{beamercolorbox}%
    \begin{beamercolorbox}[wd=.34\paperwidth,ht=2.5ex,dp=1ex,center]{date in head/foot}%
      \scriptsize 30 August 2026
    \end{beamercolorbox}%
    \begin{beamercolorbox}[wd=.33\paperwidth,ht=2.5ex,dp=1ex,right]{date in head/foot}%
      \insertframenumber{} / \inserttotalframenumber
      \hspace*{2mm}
    \end{beamercolorbox}%
  }%
}

% Remove navigation symbols
\setbeamertemplate{navigation symbols}{}

\title{\small PBADT504 BIG DATA ANALYTICS PROJECT}
\subtitle{\Huge Real-Time Social Media Anomaly Detection \\ \large Using Kafka, Spark Streaming, NLP \& Power BI}
\author{
11 Aleena Anil (CEC24AD012) \\
23 Edwin E Binu (CEC24AD024) \\
38 Niranjana TM (CEC24AD039) \\
14 Angel Sudhish (CEC24AD015) \\
25 Farhana A (CEC24AD026)
}
\institute{College of Engineering Cherthala}
\date{Academic Year 2026--27}

\begin{document}

% 1. Title
\frame{\titlepage}

% 2. Challenge
\begin{frame}{The Challenge}
\begin{block}{Problem Statement}
Social media platforms generate millions of posts every minute.
Detecting unusual trends, fake news, hate speech, or sudden spikes manually is impossible.
\end{block}

\begin{itemize}
    \item 1.6M+ tweets analyzed in dataset
    \item Validated real-time detection at scale
\end{itemize}
\end{frame}

% 3. Why It Matters
\begin{frame}{Why It Matters}
\begin{enumerate}
    \item Early detection of viral events
    \item Identify misinformation
    \item Detect abnormal user behavior
    \item Support governments and organizations
\end{enumerate}
\end{frame}

% 4. Objectives
\begin{frame}{Objectives}
\textbf{Main Objective:} Develop a real-time anomaly detection system for social media streams using Big Data technologies.

\begin{itemize}
    \item Collect: Stream data using Apache Kafka
    \item Process: Handle tweets using Apache Spark Streaming
    \item Analyze: Perform sentiment analysis using NLP
    \item Detect: Identify anomalies using ML
    \item Visualize: Present insights using Power BI
\end{itemize}
\end{frame}

% Literature Survey
\begin{frame}{Literature Survey}
\scriptsize
\begin{tabular}{|l|l|l|l|}
\hline
Paper (Author) & Year & Method & Limitation \\
\hline
C. Anitha & 2025 & LightGBM vs SVM comparative study & Limited to specific models \\
\hline
P. Upadhyay et al. & 2023 & ML-based sentiment analysis & Generalization across platforms \\
\hline
R. Bharal, O.V. Krishna & 2021 & CNN-BiLSTM deep learning & Dataset constraints \\
\hline
H. Pal, B. Bhushan & 2024 & Voting classifier on Twitter data & Class imbalance issues \\
\hline
Y. Sahu, V. Sharma, A. Bhargava & 2025 & Hybrid BiLSTM-CNN architecture & Computational complexity \\
\hline
\end{tabular}
\end{frame}


% 6. Methodology
\begin{frame}{Proposed Methodology}
\begin{enumerate}
    \item Dataset (Twitter/X Data)
    \item Apache Kafka
    \item Apache Spark Streaming
    \item Text Preprocessing
    \item Feature Extraction (TF-IDF)
    \item Machine Learning (Isolation Forest)
    \item Anomaly Detection
    \item Power BI Dashboard
\end{enumerate}
\end{frame}

% 7. Tools & Technologies
\begin{frame}{Tools \& Technologies}
\begin{columns}
\column{0.5\textwidth}
\begin{itemize}
    \item Python
    \item Apache Kafka
    \item Apache Spark
    \item Hadoop HDFS
    \item Jupyter Notebook
\end{itemize}
\column{0.5\textwidth}
\begin{itemize}
    \item VS Code
    \item NLTK, TextBlob
    \item Scikit-learn
    \item Pandas
    \item Power BI
    \item Git \& GitHub
\end{itemize}
\end{columns}
\end{frame}

% 8. Dataset
\begin{frame}{Dataset \& Preprocessing}
\textbf{Source:} Kaggle Twitter Sentiment Dataset (1.6M tweets)

\begin{itemize}
    \item Fields: Text, Timestamp, Username, Sentiment Label, Tweet ID
    \item Preprocessing:
    \begin{enumerate}
        \item Remove URLs
        \item Remove Stop Words
        \item Tokenization
        \item Stemming
        \item Lemmatization
    \end{enumerate}
\end{itemize}
\end{frame}

% 9. Work Completed
\begin{frame}{Work Completed}
\begin{itemize}
    \item Literature Survey \done
    \item Dataset Selected \done
    \item Kafka Installed \& Configured \done
    \item Spark Installed \& Streaming Job Deployed \done
    \item Hadoop Configured \done
    \item Python Environment Ready \done
    \item End-to-End Data Pipeline Built (Kafka $\rightarrow$ Spark $\rightarrow$ CSV) \done
    \item NLP Sentiment Analysis Integrated (TextBlob) \done
    \item Power BI Dashboard Built \& Connected to Live Output \done
\end{itemize}
\end{frame}

% 10. Final Results Snapshot
\begin{frame}{Final Results Snapshot}
\begin{itemize}
    \item Kafka successfully streams tweets \done
    \item Spark processes streaming data in real time \done
    \item Sentiment analysis (polarity, subjectivity, label) working end-to-end \done
    \item Power BI dashboard live and connected to pipeline output \done
    \item Brand-level and time-series sentiment breakdown delivered \done
\end{itemize}

\vspace{0.3cm}
\textbf{Pipeline Throughput:} 30 tweets processed in the validation run, sub-5-second micro-batch latency, zero dropped messages.
\end{frame}

% 11. Results — Dashboard Overview
\begin{frame}{Results: Sentiment Dashboard Overview}
\begin{center}
\includegraphics[width=0.92\textwidth]{dashboard_overview.png}
\end{center}
\end{frame}

% 12. Results — Brand Breakdown
\begin{frame}{Results: Brand-Wise Sentiment Breakdown}
\begin{center}
\includegraphics[width=0.92\textwidth]{dashboard_brand_breakdown.png}
\end{center}
\end{frame}

% 13. Key Metrics
\begin{frame}{Key Metrics}
\begin{columns}
\column{0.5\textwidth}
\textbf{Overall Sentiment Distribution}
\begin{itemize}
    \item Total tweets analyzed: \textbf{30}
    \item Positive: \textbf{13 (43.3\%)}
    \item Negative: \textbf{10 (33.3\%)}
    \item Neutral: \textbf{7 (23.3\%)}
    \item Average Polarity: \textbf{0.09}
    \item Average Subjectivity: \textbf{0.41}
\end{itemize}
\column{0.5\textwidth}
\textbf{Brand-Level Insights}
\begin{itemize}
    \item Brands tracked: \textbf{14} (Amazon, Apple, Tesla, Netflix, Microsoft, PlayStation, Xbox, Google, Spotify, Facebook, Borderlands, Call of Duty, Verizon, Other)
    \item Top mentioned category: \textbf{``Other'' (5 mentions)}
    \item Second most mentioned: \textbf{Tesla (3 mentions)}
    \item Sentiment varies meaningfully across brands, confirming entity-level signal in the data
\end{itemize}
\end{columns}
\end{frame}

% 14. Team Contributions
\begin{frame}{Team Contributions}
\begin{tabular}{|l|l|}
\hline
Niranjana TM & Literature Survey \\
\hline
Angel Sudhish & Dataset Collection \& Preprocessing \\
\hline
Edwin E Binu & Kafka + Spark Development, Sentiment Pipeline, Dashboard \\
\hline
Aleena Anil & Dashboard Development \\
\hline
Farhana A & Documentation \\
\hline
\end{tabular}
\end{frame}

% 15. Challenges
\begin{frame}{Challenges \& Solutions}
\begin{columns}
\column{0.5\textwidth}
\textbf{Challenges}
\begin{itemize}
    \item Large volume of streaming data
    \item Kafka configuration issues (KRaft mode setup on Windows)
    \item Spark memory management
    \item Missing tweet values
    \item Class imbalance
    \item Power BI folder-connector and OneDrive path conflicts
    \item Hardware limitations
\end{itemize}
\column{0.5\textwidth}
\textbf{Solutions}
\begin{itemize}
    \item Data cleaning \& validation
    \item Memory optimization (micro-batch tuning)
    \item Feature engineering
    \item Parameter tuning
    \item Switched to single-file CSV load in Power BI to bypass connector bug
\end{itemize}
\end{columns}
\end{frame}

% 16. Future Work
\begin{frame}{Future Work}
\begin{itemize}
    \item Improve anomaly detection accuracy with dedicated ML models
    \item Deploy on cloud (multi-node Spark cluster)
    \item Real-time Twitter/X API integration (replacing CSV-driven simulation)
    \item Deep Learning models (transformer-based sentiment, e.g. BERT)
    \item Live, auto-refreshing Power BI dashboard via streaming dataset
    \item Mobile notification system for sentiment spikes
\end{itemize}
\end{frame}

% 17. Timeline
\begin{frame}{Project Timeline}
\centering
\begin{tabular}{|l|c|}
\hline
\textbf{Activity} & \textbf{Status} \\
\hline
Literature Survey & \done Completed \\
Dataset Collection & \done Completed \\
Environment Setup & \done Completed \\
Kafka Integration & \done Completed \\
Spark Processing & \done Completed \\
NLP Sentiment Analysis & \done Completed \\
Power BI Dashboard & \done Completed \\
Final Documentation & \done Completed \\
\hline
\end{tabular}
\end{frame}

% 18. Final Slide
\begin{frame}
\centering
\Huge Thank You
\end{frame}

\end{document}
